% Generated draft: requires mathematical and source review.
\documentclass[11pt]{article}
\usepackage[margin=25mm]{geometry}
\usepackage{fontspec}
\setmainfont{Latin Modern Roman}
\usepackage{amsmath,amssymb,mathtools,unicode-math}
\setmathfont{Latin Modern Math}
\usepackage{enumitem,needspace,xurl,hyperref}
\setlength{\parindent}{0pt}
\setlength{\parskip}{5pt}
\setlength{\emergencystretch}{3em}
\newcommand{\N}{\mathbb{N}}
\newcommand{\Z}{\mathbb{Z}}
\newcommand{\Q}{\mathbb{Q}}
\newcommand{\R}{\mathbb{R}}
\newcommand{\C}{\mathbb{C}}
\newcommand{\F}{\mathbb{F}}
\newcommand{\A}{\mathbb{A}}
\newcommand{\PP}{\mathbb{P}}
\newcommand{\E}{\mathbb{E}}
\newcommand{\eps}{\varepsilon}
\newcommand{\dd}{\,\mathrm{d}}
\begin{document}
\section*{UnsolvedMath ID: 700314}
\textbf{Source problem code:} catalogue entry
\subsection{Definitions and mathematical statement}
Let \(G\) be a finite soluble group, meaning that its derived series terminates at the trivial group. Define its nilpotent length \(h(G)\) to be the least integer \(h\geq 0\) for which there exists a series
\[
1=G_0\leq G_1\leq\cdots\leq G_h=G
\]
with every \(G_i\) normal in \(G\) and every factor \(G_i/G_{i-1}\) nilpotent. Here a group is nilpotent if it has a finite central series, and \(h(1)=0\).

Let \(A\leq\operatorname{Aut}(G)\), and define the common fixed-point subgroup by
\[
C_G(A)=\{g\in G: a(g)=g\text{ for every }a\in A\}.
\]
Assume that
\[
C_G(A)=1,\qquad \gcd(|G|,|A|)=1.
\]
Write
\[
|A|=\prod_{i=1}^{r}p_i^{e_i},\qquad n=\sum_{i=1}^{r}e_i,
\]
where the \(p_i\) are distinct primes and each \(e_i\geq 1\). Thus \(n\) counts prime factors with multiplicity; for \(|A|=1\), the empty-product convention gives \(n=0\).

\textbf{Question.} For every such pair \((G,A)\), must
\[
h(G)\leq n
\]
hold? No solubility assumption on \(A\) is imposed in this question. In the boundary case \(A=1\), the fixed-point hypothesis forces \(G=1\).

\textbf{Supplied qualifications.} The source calls this a well-known problem and reports that the proposed bound has been proved for large classes of groups, crediting E. Shult, F. Gross, T. Berger, and A. Turull without specifying those classes. Under the additional hypothesis that \(A\) is soluble, it reports bounds \(h(G)\leq 5^n\), \(h(G)\leq 4n\), and \(h(G)\leq 2n\), attributed respectively to J. G. Thompson, H. Kurzweil, and A. Turull. These are supplied historical assertions whose precise theorem hypotheses and references are not included. The supplied record labels the problem open; this draft does not independently verify that status.
\subsection{Short English statement}
V. D. Mazurov's problem asks whether a finite soluble group admitting a coprime automorphism group with trivial common fixed-point subgroup has nilpotent length at most the number of prime factors of the automorphism group's order, counted with multiplicity. The automorphism group is not assumed soluble. The supplied source reports partial results and labels the general question open.
\subsection{Sources}
\begin{itemize}
\item Local UnsolvedMath registry, permanent ID \(700314\), exact source problem code \texttt{catalogue entry}, title \emph{Kourovka Notebook Problem 5.30}; supplied complete source record. The record carries status \texttt{open} and update timestamp \texttt{2026-10-16T14:00:00Z}; that timestamp is not identified as a literature-verification date.
\item \emph{The Kourovka Notebook (Unsolved Problems in Group Theory)}, original Problem 5.30, fifth issue (1976), proposed by V. D. Mazurov. The supplied source identifies the twenty-first issue edition (2026), \texttt{arXiv:1401.0300v48}; its collection description names E. I. Khukhro and V. D. Mazurov as editors. This attribution is reproduced from the record without independent source consultation.
\item The supplied statement credits E. Shult, F. Gross, T. Berger, and A. Turull for results on large classes, and J. G. Thompson, H. Kurzweil, and A. Turull for the respective bounds \(5^n\), \(4n\), and \(2n\) when \(A\) is soluble. No individual publication references or theorem statements are supplied for these results.
\end{itemize}
\end{document}
